\section{Implementación y Resultados Técnicos}\label{cap:implementacion}

Este capítulo documenta lo que existe construido y verificado al 26 de noviembre de 2026 (semana 4 del cronograma). Todas las cifras proceden de ejecuciones realizadas sobre el repositorio en esa fecha o del conteo de su código y esquema; no provienen de datos reales del salón. Los resultados de efecto sobre el tiempo de liquidación, los errores de cálculo y la usabilidad (Capítulo~\ref{cap:operacionalizacion}) \textbf{no se reportan aquí}, porque requieren la medición posterior descrita en el marco metodológico.

\subsection{Alcance construido}

El sistema se implementó como un monolito modular: cada dominio tiene su propia carpeta con reglas puras (\textit{domain}), casos de uso (\textit{services}), adaptador HTTP (\textit{api}), persistencia y un contrato público, de modo que un módulo pueda extraerse como microservicio sin reescribirlo. La Tabla~\ref{tab:impl_1} resume los módulos de backend en operación.

\vspace{8pt}{\small\setstretch{1.0}
\begin{longtable}{>{\raggedright\arraybackslash}p{0.17\linewidth}>{\raggedright\arraybackslash}p{0.60\linewidth}>{\centering\arraybackslash}p{0.15\linewidth}}
\caption{Módulos de backend implementados}\label{tab:impl_1}\\
\toprule
\textbf{Módulo} & \textbf{Responsabilidad} & \textbf{Métodos de prueba} \\
\midrule\endfirsthead
\toprule
\textbf{Módulo} & \textbf{Responsabilidad} & \textbf{Métodos de prueba} \\
\midrule\endhead
\bottomrule\endfoot
\bottomrule\endlastfoot
identity & Salones, usuarios, roles, permisos y autenticación & 28 \\
audit & Historial de cambios con contexto de usuario & 4 \\
catalog & Servicios, especialistas y clientas & 9 \\
commissions & Reglas de comisión versionadas y snapshot por línea de venta & 34 \\
sales & Ventas multilínea, pagos divididos, anulación e idempotencia por UUID & 53 \\
cashier & Resumen y cierre de caja por método de pago & 30 \\
payroll & Deducciones, liquidación semanal, aprobación, pago, PDF y XLSX & 79 \\
analytics & Panel financiero (lectura) & 19 \\
notifications & Bandeja de salida (\textit{outbox}) para n8n & 28 \\
demo & Datos sintéticos para la demostración & 7 \\
\end{longtable}}
\nota{Elaboración propia a partir de \texttt{Backend/lushbeauty/apps/}. Se contaron los métodos \texttt{test\_*} de cada módulo (291); el ejecutor de pruebas reporta 295 casos porque algunos métodos generan más de un caso.}

En el frente de usuario existen dos clientes: la aplicación web progresiva (PWA) de React y TypeScript, organizada por módulos (\textit{sales}, \textit{payroll}, \textit{cashier}, \textit{analytics}, \textit{decisions}, \textit{insight}, \textit{guard}, \textit{imports}, entre otros), y una réplica móvil en Flutter con las vistas de la especialista (ingreso, \textit{Hoy}, registro rápido, sincronización, \textit{Mi liquidación} y cierre de caja). La base de datos PostgreSQL se define en un esquema propio con 23 tablas y 20 políticas de seguridad a nivel de fila (RLS).

\subsection{Verificación ejecutada}

La Tabla~\ref{tab:impl_2} presenta las suites ejecutadas el 26 de noviembre de 2026. El backend se probó contra PostgreSQL dentro de Docker; las pruebas de dominio (comisión, contribución, liquidación) usan \texttt{Decimal} y no requieren base de datos.

\vspace{8pt}{\small\setstretch{1.0}
\begin{longtable}{>{\raggedright\arraybackslash}p{0.28\linewidth}>{\raggedright\arraybackslash}p{0.28\linewidth}>{\centering\arraybackslash}p{0.14\linewidth}>{\raggedright\arraybackslash}p{0.20\linewidth}}
\caption{Pruebas automatizadas ejecutadas}\label{tab:impl_2}\\
\toprule
\textbf{Componente} & \textbf{Herramienta} & \textbf{Casos} & \textbf{Resultado} \\
\midrule\endfirsthead
\toprule
\textbf{Componente} & \textbf{Herramienta} & \textbf{Casos} & \textbf{Resultado} \\
\midrule\endhead
\bottomrule\endfoot
\bottomrule\endlastfoot
Backend Django/DRF & \texttt{manage.py test}, PostgreSQL 17 & 295 & Todos aprobados \\
Web (PWA) & Vitest & 120 & Todos aprobados \\
Móvil (Flutter) & \texttt{flutter test} & 32 & Todos aprobados \\
\end{longtable}}
\nota{Elaboración propia. El número de casos de cada suite es el informado por su ejecutor. No se midió cobertura de código; por tanto, estas cifras indican qué comportamientos están cubiertos, no qué porcentaje del código lo está.}

Entre los comportamientos comprobados por las pruebas se encuentran los siguientes casos de aceptación del Anexo~J:

\begin{vinetas}
\item \textbf{CP-01, descuento sin alterar la lista.} Un servicio de lista Bs~120 con descuento de Bs~8 se cobra a Bs~112; se conserva el precio de lista y la comisión (35\,\% ilustrativo) se calcula sobre la lista: Bs~42,00. Se verifica en dominio, API y pantalla.
\item \textbf{Idempotencia.} Reenviar una venta con el mismo UUID devuelve la existente (201 la primera vez, 200 después) sin crear filas; el mismo UUID en otro salón es otra venta.
\item \textbf{Integridad monetaria.} Pagos distintos de las líneas se rechazan sin dejar filas; el umbral de descuento de 10\,\% pasa sin aprobación y un centavo más la exige.
\item \textbf{Inmutabilidad.} La base de datos no permite borrar ni modificar líneas confirmadas ni liquidaciones aprobadas; la anulación conserva las líneas y queda auditada.
\item \textbf{Aislamiento por salón y por rol.} Un salón no ve datos de otro; una especialista solo ve sus servicios y únicamente sus liquidaciones aprobadas, con su comprobante PDF.
\end{vinetas}

\subsection{Ciclo operativo demostrable}

El ciclo central de caja y nómina funciona de extremo a extremo en el backend: registro de venta con pagos $\rightarrow$ cálculo de comisión por \textit{snapshot} versionado $\rightarrow$ cierre de caja por método $\rightarrow$ borrador y aprobación de la liquidación semanal (sueldo fijo, comisión, propina y deducciones con motivo) $\rightarrow$ marcado como pagada y emisión de comprobante PDF y resumen XLSX. La fórmula de liquidación se aplica igual en el dominio y en una restricción de la base de datos. Los eventos de cierre y de liquidación se escriben en la bandeja de salida para su entrega por n8n.

\subsection{Estado por funcionalidad}

La Tabla~\ref{tab:impl_3} distingue lo implementado de lo que sigue como maqueta. Se toma del backlog del repositorio (\texttt{odd/tasks/INDEX.md}) y se contrastó con el código y las pruebas.

\vspace{8pt}{\small\setstretch{1.0}
\begin{longtable}{>{\raggedright\arraybackslash}p{0.30\linewidth}>{\raggedright\arraybackslash}p{0.34\linewidth}>{\raggedright\arraybackslash}p{0.28\linewidth}}
\caption{Estado de las funcionalidades al 26 de noviembre de 2026}\label{tab:impl_3}\\
\toprule
\textbf{Funcionalidad} & \textbf{Backend} & \textbf{Interfaz} \\
\midrule\endfirsthead
\toprule
\textbf{Funcionalidad} & \textbf{Backend} & \textbf{Interfaz} \\
\midrule\endhead
\bottomrule\endfoot
\bottomrule\endlastfoot
Base y seguridad (F00, F14) & Hecho; RLS, auditoría y datos demo; parcial en despliegue y CI & Solo el inicio de sesión consume la API \\
Caja, ventas y pagos (F01) & Hecho & Maqueta sin conectar \\
Comisiones versionadas (F02) & Hecho & Lectura con datos de prueba \\
Liquidación semanal (F03) & Hecho, con PDF y XLSX & Maqueta \\
Cierre de caja (F04) & Hecho & Maqueta \\
PWA sin conexión (F05) & Idempotencia lista; cola y sincronizador pendientes & Maqueta \\
Importador de Excel (F06) & Sin empezar & Maqueta \\
Panel financiero (F07) & Lectura hecha & Maqueta \\
Contribución y What-If (F08) & Parcial & Maqueta \\
Matriz y sensibilidad (F09) & Sin empezar & Maqueta \\
Recomendación y decisiones (F10) & Sin empezar & Maqueta \\
Lush Insight (F11) y Guard (F12) & Sin empezar & Maqueta / respuestas fijas \\
n8n y \textit{outbox} (F13) & Eventos de cierre y liquidación; faltan los flujos & Sin empezar \\
\end{longtable}}
\nota{Elaboración propia. «Maqueta» significa pantalla con datos en memoria que no llama al API correspondiente, aunque el endpoint exista.}

\subsection{Decisiones abiertas con valor provisional}

Algunas reglas se aplicaron con un valor inicial a la espera de validación con la administradora: el redondeo de la comisión (mitad hacia arriba, a dos decimales), el umbral de descuento sin aprobación (10\,\%), los métodos de pago (efectivo, QR, transferencia y tarjeta), la semana de liquidación (lunes a domingo, configurable por salón) y los pesos por defecto de la matriz de decisión (40/25/20/15, de ejemplo). Los resultados que dependen de ellos deben leerse como condicionados a esa validación.

\subsection{Limitaciones del avance}

\begin{vinetas}
\item La interfaz web aún no consume la mayoría de los endpoints existentes; la integración convertirá las maquetas en pantallas reales.
\item El módulo de soporte para decisiones (escenarios, matriz, sensibilidad, recomendación y seguimiento) carece de backend, salvo un avance parcial del modelo de contribución.
\item No se han hecho pruebas de carga, de concurrencia real ni de campo con conexión intermitente.
\item Las capturas de pantalla de las vistas de la PWA y de Flutter se incorporarán cuando estén conectadas al API; hasta entonces no se presentan como evidencia de funcionamiento.
\item No hay mediciones de tiempo de liquidación, tasa de errores ni usabilidad (SUS); las hipótesis siguen sin contrastarse.
\end{vinetas}
